\chapter{Concurrency in object-oriented Perl applications}
\label{ch:enterprise-objects}

A production Perl application spans many lifetimes: interpreter, worker,
request, transaction and durable service. In a synthetic inventory application,
a web controller, an import job and a notification worker may all operate on
the same record. A lock attached to a Moose object cannot coordinate
independent processes. Choose the owner of the invariant first.

\section{State ownership and object lifetime}
\begin{figure}[!htbp]\centering
\begin{tikzpicture}[node distance=6mm and 12mm]
\node[box,text width=6cm] (p) {Parent: immutable configuration and class metadata};
\node[okbox,below left=of p,text width=3.5cm] (a) {Worker A\\PID-owned clients};
\node[okbox,below right=of p,text width=3.5cm] (b) {Worker B\\PID-owned clients};
\draw[arr](p)--(a); \draw[arr](p)--(b);
\node[box,below=of a,text width=3.5cm] (ra) {Request context A};
\node[box,below=of b,text width=3.5cm] (rb) {Request context B};
\draw[arr](a)--(ra); \draw[arr](b)--(rb);
\node[okbox,below=8mm of $(ra.south)!0.5!(rb.south)$,text width=6cm] (db) {Shared authoritative database};
\draw[arr](ra.south)--(db.north west); \draw[arr](rb.south)--(db.north east);
\end{tikzpicture}
\caption{Each process owns its clients; request contexts are transient and shared invariants are enforced at the resource.}
\end{figure}

\begin{longtable}{L{.22\linewidth}L{.33\linewidth}L{.37\linewidth}}
\toprule Lifetime & Appropriate state & Avoid \\ \midrule\endhead
Application & Immutable configuration, class metadata & Current user, mutable totals \\
Worker & Connection factories, bounded pools, metrics & Assuming other workers see local changes \\
Request / job & Identity, operation ID, validated input & Storing request context on a reusable singleton \\
Durable service & Inventory, claims, checkpoints, audit rows & Treating a cache as authoritative \\
\bottomrule
\end{longtable}

Moose lazy attributes defer construction; they are not synchronization
primitives [63]. Catalyst component factories and context-aware lookup
(\code{COMPONENT} and \code{ACCEPT_CONTEXT}) allow different lifetimes [64].
Pass request-specific values explicitly instead of storing the active user
or current transaction on a persistent controller. A process-local singleton
does not become a deployment-wide singleton. Interpreter threads clone ordinary
state; an event-loop task can race another task across an \code{await} even
without parallel CPU execution.

\section{Resource ownership after fork}
Prefer creating live clients after the server forks. Use documented
driver-specific fork handling and reconnect behavior. A PID check detects
inherited state, but cannot make an inherited socket safe to use or destroy.
DBI's \code{AutoInactiveDestroy} protects against accidental destructor
interference in a forked child [15]. It does not make a shared connection safe.

\warn{Never use an inherited database handle concurrently in a parent and
child. Do not call a protocol-level disconnect on an inherited client
unless its documentation establishes that this is safe.}

The following Moose laboratory uses a plain hash to model resource ownership,
not a real database connection. The child must initialize its own resource
and leave the parent unchanged.
\CodeFile{perl}{examples/32-worker-lifecycle.pl}{Example 32: process-local Moose resource ownership}

\section{Reusable services and invariants}
Give each service operation explicit inputs, a stable operation ID and
well-defined results: success, conflict, validation failure or unknown
outcome. Web, batch and administrative entry points should use the same
business rules. Prefer a resource-enforced guard or unique constraint so
that an alternate writer cannot silently bypass an advisory lock.

\chapter{Transaction boundaries in layered applications}
\label{ch:enterprise-transactions}
A service operation should define the atomic unit of work. The controller
validates the request; repository methods participate in the transaction
rather than committing independently. A database commit belongs to the
business operation, not to an individual ORM result object.

\begin{figure}[!htbp]\centering
\begin{tikzpicture}[node distance=5mm]
\node[box,text width=10cm] (c) {Controller or job handler: validate and assign operation ID};
\node[okbox,below=of c,text width=10cm] (s) {Service: begin transaction; define retry boundary};
\node[box,below left=8mm and 0mm of s.south,text width=4.2cm,anchor=north east] (a) {Repository A\\guarded stock update};
\node[box,below right=8mm and 0mm of s.south,text width=4.2cm,anchor=north west] (b) {Repository B\\audit or outbox insert};
\node[okbox,below=8mm of $(a.south)!0.5!(b.south)$,text width=10cm] (d) {Commit together or roll back together};
\draw[arr](c)--(s); \draw[arr](s)--(a); \draw[arr](s)--(b);
\draw[arr](a)--(d); \draw[arr](b)--(d);
\node[box,below=of d,text width=10cm] (r) {Return response; relay external effects separately};
\draw[arr](d)--(r);
\end{tikzpicture}
\caption{A layered application has one explicit business transaction boundary.}
\end{figure}

\section{DBIx::Class composition}
Use \code{txn_do} for a complete unit of work or \code{txn_scope_guard}
when a lexical transaction guard better expresses the lifecycle.
Nested \code{txn_do} calls do not independently commit: the outer transaction
owns the commit [65]. Savepoints, when supported and correctly configured,
allow partial rollback, not independent durability. Do not catch a failed
nested operation and continue as if the transaction were still valid.
Avoid mixing raw DBI transaction calls with ORM-managed transaction depth.

\section{Retry the operation, not the final statement}
InnoDB deadlocks roll back a transaction; a lock-wait timeout normally
rolls back only the failed statement [67]. Explicitly abandon the failed
unit before retrying it. Capture driver error details promptly, use a
bounded retry budget with jitter, and reacquire all data needed by the
operation. A connection failure during commit can leave the outcome
unknown. Resolve it by a durable request ID rather than blindly retrying
a charge, decrement or external notification.

\begin{longtable}{L{.30\linewidth}L{.62\linewidth}}
\toprule Failure & Service behavior \\ \midrule\endhead
Validation error & Return a business error; do not retry. \\
Deadlock & Roll back/reset the unit; retry from the beginning with bounded backoff. \\
Lock wait timeout & Roll back the entire business operation explicitly before retry. \\
Unique-key conflict & Determine whether this is a duplicate operation or a different payload. \\
Lost commit response & Reconcile by durable operation ID; do not assume rollback. \\
\bottomrule
\end{longtable}

\section{Laboratory: a reservation and its audit record}
The example below deliberately fails between a stock decrement and an audit
insert, verifies rollback, then verifies successful commit and duplicate
rejection. SQLite is used only to make this application-layer test
self-contained; the MySQL chapter covers InnoDB-specific lock semantics.
\CodeFile{perl}{examples/33-layered-transaction.pl}{Example 33: service-owned DBIx::Class transaction}

\tip{Keep expensive computation and network I/O outside the transaction.
If the computation depends on mutable database state, validate its version
again inside the transaction before committing.}

\chapter{Scheduled jobs and resumable processing pipelines}
\label{ch:enterprise-pipelines}
Consider an import of measurement files or transaction records, followed by
validation, enrichment, aggregation and publication. The input can arrive
twice; a worker can die halfway through a file; a downstream service can
temporarily fail. The design must preserve completeness and avoid duplicate
business effects without assuming that a job executes exactly once.

\section{Pipeline decomposition}
\begin{figure}[!htbp]\centering
\begin{tikzpicture}[node distance=5mm]
\node[box,text width=2.3cm] (i) {Input\\manifest};
\node[box,right=of i,text width=2.3cm] (v) {Validate\\and stage};
\node[okbox,right=of v,text width=2.3cm] (q) {Durable\\job queue};
\node[box,right=of q,text width=2.3cm] (o) {Publish\\results};
\draw[arr](i)--(v); \draw[arr](v)--(q);
\node[badbox,below=of v,text width=2.3cm] (e) {Quarantine\\invalid input};
\node[okbox,below=of q,text width=2.3cm] (w) {Bounded\\worker pool};
\draw[badarr](v)--(e); \draw[arr](q)--(w); \draw[arr](w.east)-|(o.south);
\node[lbl,below=of w,text width=5cm] {Claim, process, checkpoint; retry with the same input identity};
\end{tikzpicture}
\caption{A durable processing pipeline separates admission, execution and publication.}
\end{figure}
Assign a stable source identity such as a producer ID plus record ID, or
a file manifest ID plus record offset. Validate the identity and payload
together. Stage data in chunks, quarantine invalid records, and use a unique
constraint on the canonical input identity. If a replay changes the payload
for the same identity, surface a conflict rather than silently overwrite it.

Partition work by a key that owns the invariant: account, device, customer
or dataset. Independent partitions can run concurrently; dependent updates
must be ordered or guarded. Use bounded queues and cap the number of workers
by the downstream connection and write budget. A checkpoint is valid only
when the effects it describes are durable. A high-water mark alone is
unsafe when records finish out of order: track completed ranges or individual
identities until the contiguous prefix is complete.

\section{Prevent overlapping scheduled runs}
Host-local coordination is appropriate when all launchers execute on the
same host and share a stable lock file. The Linux util-linux \code{flock}
command can return a distinct conflict code; preserve it in monitoring.
\CodeFile{bash}{snippets/enterprise-cron.sh}{Host-local cron wrapper (deployment excerpt)}
Do not delete or rotate the lock file while a process may hold it.
Do not let the protected command detach or close the inherited lock
descriptor. A missed run should be observable; do not confuse it with a
successful import. Across hosts use durable claims or a suitable shared
coordinator rather than unrelated local lock files.

\section{Durable claims, attempts and deadlines}
\begin{figure}[!htbp]\centering
\begin{tikzpicture}[node distance=9mm]
\node[box,text width=2.1cm] (r) {Ready};
\node[okbox,right=of r,text width=2.5cm] (w) {Running\\attempt $n$};
\node[okbox,right=of w,text width=2.1cm] (d) {Done};
\node[badbox,below=of w,text width=2.5cm] (x) {Dead letter};
\draw[arr](r)--(w) node[midway,lbl,above]{claim};
\draw[arr](w)--(d) node[midway,lbl,above]{complete};
\draw[badarr](w.south west) to[bend left=20] node[lbl,below left]{expire / retry} (r.south);
\draw[badarr](w)--(x) node[midway,lbl,right]{budget spent};
\end{tikzpicture}
\caption{A claim identifies an attempt, not permanent ownership of a job.}
\end{figure}
A worker claims one row in a short transaction, performs expensive work
outside it, then conditionally completes the matching attempt. A reclaimed
job has a newer attempt number. The earlier worker must not be allowed
to mark the replacement's work complete. Enforce uniqueness or fencing at
the effect destination as well; changing job status alone cannot prevent
a paused worker from executing a side effect.

\CodeFile{sql}{snippets/outbox-claim.sql}{MySQL claim and conditional completion (SQL sequence)}
The placeholders represent bound values from the selected row. Use an
InnoDB job table with a primary key, status, attempt, lease expiry and
indexes appropriate to the claim predicates; a production scheduler also
needs a retry budget, next eligible time and dead-letter handling.
Use \code{SKIP LOCKED} for queue-like access, not for general consistent
reads. Check each conditional update's affected-row count. If no row is
selected, roll back and wait using a bounded polling interval.

\section{Operational examples}
\begin{longtable}{L{.26\linewidth}L{.34\linewidth}L{.32\linewidth}}
\toprule Scenario & Concurrency boundary & Replay contract \\ \midrule\endhead
File import & Manifest and input identity & Reject divergent duplicate payloads \\
Aggregation & Partitioned dataset and period & Recompute or atomically replace a version \\
Report generation & Consistent read snapshot & Publish completed output by version \\
Export & Durable delivery job & Retry using stable destination identity \\
\bottomrule
\end{longtable}
For scheduled jobs on Kubernetes, \code{concurrencyPolicy: Forbid} prevents
overlap only among Jobs belonging to the same CronJob [66]. It is not a
distributed business lock, and scheduling is approximate. Jobs must remain
idempotent even when manually started or recreated.
\CodeFile{yaml}{snippets/enterprise-cronjob.yaml}{Kubernetes scheduling policy (illustrative manifest)}

\chapter{Reliable cross-system integration}
\label{ch:enterprise-outbox}
A local SQL transaction cannot roll back an HTTP request or a message
already accepted by another system. Holding a local transaction open while
calling the external system does not solve this problem; it also increases
lock duration. Instead, separate durable intent from delivery.

\section{Transactional outbox and inbox}
\begin{figure}[!htbp]\centering
\begin{tikzpicture}[node distance=10mm]
\node[okbox,text width=5cm] (s) {Source database\\business row + outbox event\\one transaction};
\node[box,right=of s,text width=3cm] (r) {Relay\\claim and send};
\node[okbox,below=of r,text width=3cm] (d) {Destination\\unique event ID};
\node[box,left=of d,text width=5cm] (a) {Source acknowledgement\\conditional on claim};
\draw[arr](s)--(r); \draw[arr](r)--(d) node[midway,lbl,right]{stable ID};
\draw[arr](d)--(a); \draw[arr](a)--(s);
\node[badbox,below=of a,text width=5cm] (f) {Crash after remote commit?\\Retry with the same ID};
\draw[badarr](f.east)-|(d.south);
\end{tikzpicture}
\caption{The outbox preserves intent; destination deduplication resolves ambiguous delivery.}
\end{figure}
Insert an outbox event alongside the business row in the same transaction
[71]. A relay claims committed events, delivers them, and records completion
after acknowledgement. If the response is lost, the relay cannot know
whether the remote operation happened. Reuse the event ID on retries.
The receiving application should commit its deduplication record (inbox)
and effect in a single transaction. A unique inbox row inserted before a
nontransactional external effect does not, by itself, guarantee exactly
one effect.

\section{Failure windows and recovery}
\begin{longtable}{L{.34\linewidth}L{.58\linewidth}}
\toprule Failure point & Recovery \\ \midrule\endhead
Before source commit & Neither business state nor intent is committed; replay input. \\
After source commit & Outbox remains pending; relay later. \\
After destination commit, before acknowledgement & Retry same event ID; destination returns prior result or deduplicates. \\
After acknowledgement & No delivery needed; retain audit and dedupe records for the replay horizon. \\
\bottomrule
\end{longtable}
External APIs vary: some support idempotency keys, others offer lookup by a
client reference, and some offer neither. For the last case, ambiguous
outcomes may require reconciliation or human review. Compensation is a new
business operation, not a rollback of history. Preserve ordering by
aggregate when required, and record failures with correlation identifiers
rather than credentials or full sensitive payloads.

\section{Laboratory: two independent databases}
The following laboratory uses separate local SQLite databases, one for
business state and outbox and one for the receiving effect. It injects a
failure after the destination commits but before the source acknowledges,
then replays the same event. It also verifies source rollback and rejects
reuse of an input identity with a different value. The relay is intentionally
single-worker; multiple relays need the claim protocol of
Chapter~\ref{ch:enterprise-pipelines}.
\CodeFile{perl}{examples/34-replay-outbox.pl}{Example 34: rollback, replay and destination deduplication}

\chapter{Modernizing legacy locking and database access}
\label{ch:enterprise-migration}
Legacy applications often contain multiple generations of synchronization:
file locks around scripts, table locks around bulk writes, named advisory
locks around cross-table operations, and transactional row locks in newer
code. Removing a legacy lock without identifying every writer can create
a concurrency bug that did not previously exist.

\section{Inventory before migration}
\begin{figure}[!htbp]\centering
\begin{tikzpicture}[node distance=5mm]
\node[box,text width=10cm] (a) {Inventory writers, storage engines, locks and transaction owners};
\node[box,below=of a,text width=10cm] (b) {Measure contention and capture regression fixtures};
\node[okbox,below=of b,text width=10cm] (c) {Add guards, constraints and replay-safe identities};
\node[okbox,below=of c,text width=10cm] (d) {Canary new writers; compare and reconcile outcomes};
\node[box,below=of d,text width=10cm] (e) {Retire old locking only when every writer obeys the new contract};
\draw[arr](a)--(b); \draw[arr](b)--(c); \draw[arr](c)--(d); \draw[arr](d)--(e);
\end{tikzpicture}
\caption{Incremental migration preserves the invariant across mixed generations of code.}
\end{figure}
Inspect table engines, foreign keys, unique constraints, isolation levels,
queries, cron launchers and administrative scripts. MyISAM is nontransactional;
changing a DBI transaction boundary cannot make a MyISAM update roll back.
A move to InnoDB changes locking, index and operational requirements [68].
Test migration and rollback procedures on representative disposable data.

\section{Table locks and connection-local state}
MySQL \code{LOCK TABLES}, DDL and some administrative statements can
interact with transaction boundaries through implicit commits. Consult the
server's version-specific documentation [11,67]. Avoid combining
table-locking legacy code with ORM-managed transactions without a precise
contract. Named advisory locks, temporary tables, session variables and
isolation overrides are connection-local. A pooled connection must be reset
before reuse, or explicitly dedicated to the operation.

Long-running reports can retain snapshots and delay purge or contend with
writers, depending on the query and isolation. Prefer appropriate indexes,
bounded report windows and consistent snapshots. A replica may be useful
for reports that tolerate lag, but not for correctness-critical read-after-write
decisions unless the required consistency is enforced.

\warn{An online schema change can still require metadata locks. Test the
DDL algorithm and lock requirements for the exact server and operation;
do not assume that an online label means zero blocking.}

\section{Migration acceptance criteria}
Before removing a lock, list every code path that can mutate the protected
resource. Establish an invariant test, deploy the new guard while preserving
compatibility, and compare results with representative load. Record a
rollback strategy that does not reintroduce mixed incompatible writers.
Avoid a big-bang migration from local to distributed locking when a database
constraint can solve the problem more directly.

\chapter{Concurrency verification in application test suites}
\label{ch:enterprise-testing}
Concurrency tests should prove an invariant, not merely demonstrate that
two workers finished. Random sleeps can help stress a system but are not
a reliable method of forcing a specific race. Use rendezvous points,
independent connections, explicit deadlines and failure injection.

\section{Force the interleaving}
\begin{figure}[!htbp]\centering
\begin{tikzpicture}[x=1.35cm,y=.8cm]
\node[anchor=east] at (0,2){Worker A}; \node[anchor=east] at (0,0){Worker B};
\draw[lane](0,2)--(7.5,2); \draw[lane](0,0)--(7.5,0);
\foreach \y in {0,2}{
  \node[box] at (1,\y){Read 0};
  \node[box] at (3,\y){Ready};
  \node[badbox] at (5.8,\y){Write 1};
  \draw[arr](1.65,\y)--(2.3,\y);
  \draw[arr](3.65,\y)--(5.0,\y);
}
\draw[teal,thick](4.4,-.45)--(4.4,2.45);
\node[lbl,above] at (4.4,2.45){Release both workers};
\node[badbox] at (7.3,1){Final 1\\expected 2};
\draw[badarr](6.45,2)--(6.9,1.45);
\draw[badarr](6.45,0)--(6.9,.55);
\end{tikzpicture}
\caption{A deterministic test makes both readers observe the same old value before either writes.}
\end{figure}
The laboratory first protects individual reads and writes with locks,
which is insufficient for a compound update. A pair of queues forces
the lost-update interleaving. It then locks the complete read-modify-write
and verifies the exact final count.
\CodeFile{perl}{examples/35-deterministic-race.pl}{Example 35: reproducible lost update and corrected critical section}

\section{Parallel test isolation}
Test2::Harness can run multiple test files concurrently [69].
Parallelizing a test runner is separate from testing concurrency inside a
single test. Give parallel tests separate schemas or uniquely scoped data,
independent temporary directories and distinct service namespaces. Do not
flush a shared cache or reset a shared schema while another test is running.
Serialize tests that cannot yet be isolated, but treat that as a migration
constraint rather than a correctness guarantee.

\begin{longtable}{L{.31\linewidth}L{.61\linewidth}}
\toprule Test scenario & Invariant to assert \\ \midrule\endhead
Opposing transactions & Conserved totals; no negative balance; bounded retries \\
Worker crash after claim & Job becomes recoverable; old attempt cannot complete it \\
Lost remote response & Same event ID results in one durable effect \\
Duplicate input & Original payload preserved; divergent duplicate rejected \\
Process fork & Child uses a new connection; parent remains operational \\
Rolling restart & Admission stops; committed work remains recoverable \\
\bottomrule
\end{longtable}
Use deadlines at both the test and job level so a failed barrier cannot
hang CI. Run stress tests repeatedly with logged random seeds in addition
to deterministic tests. A test against SQLite is useful for application
composition but cannot prove MySQL-specific locking or deadlock behavior.
Run actual engine integration tests for those claims.

\chapter{Production deployment and observability}
\label{ch:enterprise-operations}
Concurrency controls must remain valid as worker counts, replicas,
schedulers and service availability change. A configuration that works on
a single development process may fail when a production deployment forks
workers or rolls out multiple replicas.

\section{Budget workers and connections}
For a hypothetical deployment of four replicas, six web workers per
replica and one live database connection per worker, web traffic alone can
consume 24 database connections. Add batch workers, administrative sessions,
pool overflow and rollout surge to the budget. A connection limit is not
a throughput target. Start with bounded concurrency and measure queueing
and database saturation before increasing parallelism.

Use the same resource ownership discipline with Apache, FastCGI and PSGI
workers: inspect the actual server lifecycle, avoid inheriting live clients,
and distinguish process-local from cluster-wide locks. Maintain bounded
request and job deadlines. Never hold an ordinary synchronization lock
across an unbounded external call.

\section{Graceful termination and durable recovery}
\begin{figure}[!htbp]\centering
\begin{tikzpicture}[node distance=7mm]
\node[box,text width=2.7cm] (t) {Termination requested};
\node[okbox,right=of t,text width=2.7cm] (s) {Stop accepting work};
\node[okbox,right=of s,text width=2.7cm] (d) {Drain bounded in-flight work};
\node[box,below=of d,text width=2.7cm] (e) {Commit, close, exit};
\node[badbox,below=of s,text width=2.7cm] (r) {Deadline expires\\reclaim later};
\draw[arr](t)--(s); \draw[arr](s)--(d); \draw[arr](d)--(e);
\draw[badarr](d.south west)--(r.north east);
\end{tikzpicture}
\caption{Graceful shutdown is an optimization; durable recovery also handles forced termination.}
\end{figure}
On a shutdown request, mark the process unready, stop claiming jobs,
finish bounded in-flight work, then close resources. A Kubernetes
termination grace period is finite [72]; forced termination must not
require a final cleanup handler for correctness. Expired attempts,
idempotent effects and durable checkpoints provide recovery when a
worker cannot finish.

\section{Measure symptoms, not just CPU}
\begin{longtable}{L{.33\linewidth}L{.59\linewidth}}
\toprule Signal & Meaning and action \\ \midrule\endhead
Lock wait duration & Histogram; increasing tail latency identifies contention. \\
Lock hold duration & Histogram; investigate slow critical sections or transactions. \\
Job queue age & Oldest eligible job; detects starvation or insufficient capacity. \\
Attempt and retry counters & Separate conflicts, transient failures and exhausted budgets. \\
Database connections & Active, idle and waiting; check saturation and leaked sessions. \\
Outbox age and backlog & Delivery delay or remote outage; alert on sustained lag. \\
\bottomrule
\end{longtable}
Trace the business operation, transaction, claim and delivery separately.
Record request ID, job ID, attempt, worker PID and elapsed time in structured
logs. Keep high-cardinality identifiers out of metric labels; use bounded
operation and outcome categories instead [73]. Never log tokens, passwords
or private payloads. Correlate traces and logs with database lock diagnostics.

\tip{Start with a small dashboard: queue age, throughput, retry rate,
lock wait percentiles, connection usage and failed jobs. Use these signals
to choose the next concurrency improvement rather than tuning blindly.}

\section{Architecture review checklist}
\begin{itemize}
\item Identify every writer and the resource that enforces its invariant.
\item Document transaction scope, lock order and retry classification.
\item Bound worker counts, queues, connection use and operation deadlines.
\item Define durable identities, replay rules and ambiguous-outcome recovery.
\item Verify crash, timeout, retry and duplicate delivery scenarios.
\item Establish deployment, rollback and observable failure criteria.
\end{itemize}
